\let\oldprintchaptertitle=\printchaptertitle
\renewcommand{\printchaptertitle}[1]{%
	\vspace*{-75pt}
	\oldprintchaptertitle{#1}
}%
\myChapterStar{R\'esum\'e}{}{section}
\let\printchaptertitle=\oldprintchaptertitle

La surveillance de la qualit\'e des eaux souterraines face aux substances per- et polyfluoroalkyl\'ees (PFAS) soul\`eve un d\'efi op\'erationnel majeur. Depuis l'adoption en 2024 par l'EPA am\'ericaine des premi\`eres limites maximales de contaminants obligatoires (jusqu'\`a 4~ng/L pour le PFOA et le PFOS dans le cadre de la r\'eglementation NPDWR), les gestionnaires de r\'eseaux doivent prioriser les campagnes d'\'echantillonnage sans pouvoir soumettre chaque forage \`a une analyse chromatographique co\^uteuse. Ce besoin d\'efinit le \emph{mode pr\'edictif strict}~: estimer le risque de d\'epassement r\'eglementaire \`a partir de variables contextuelles g\'eospatiales seules, sans concentration PFAS pr\'ealablement mesur\'ee. Ce travail adresse deux verrous que la litt\'erature existante ne traite pas simultan\'ement. Le premier est logiciel~: l'enrichissement g\'eospatial des observations brutes est g\'en\'eralement assur\'e par des scripts monolithiques non r\'eutilisables. Nous proposons \texttt{geoenrich}, un moteur g\'en\'erique organis\'e en trois couches (couche \emph{contrat} agnostique, couche \emph{op\'erateurs} param\'etrable, couche \emph{configuration} d\'eclarative), dont la correction a \'et\'e valid\'ee par un test de parit\'e avec le pipeline d'origine~: 77~colonnes identiques, 17~am\'elior\'ees et aucune r\'egression sur 94~colonnes produites. Le second verrou est scientifique et tient \`a l'auto-corr\'elation spatiale de la contamination PFAS, qui rend le d\'ecoupage al\'eatoire train-test optimiste. Pour y r\'epondre, toute variable de localisation a \'et\'e exclue des entr\'ees, et cinq mod\`eles ont \'et\'e entra\^in\'es \`a protocole identique sur 46\,338 pr\'el\`evements californiens~: deux tabulaires de r\'ef\'erence (for\^et al\'eatoire et XGBoost), un transformeur de graphe h\'et\'erog\`ene (HGT) exploitant la structure relationnelle entre forages et sources, et deux hybrides (Embedding Fusion et Stacking). Sur le jeu de test, la for\^et al\'eatoire et le stacking atteignent tous deux F1~0,925 (AUC 0,979 et 0,977)~; XGBoost obtient F1~0,923 (AUC~0,977)~; le HGT reste en retrait (F1~0,825, AUC~0,906). L'analyse par valeurs de Shapley confirme que les variables de contexte hydrog\'eologique et de proximit\'e aux sources portent l'essentiel du signal pr\'edictif, l\'egitimant le retrait des coordonn\'ees et montrant qu'un enrichissement soign\'e suffit \`a atteindre une discrimination op\'erationnelle.

\vspace{1cm}
\noindent\textbf{Mots cl\'es~:} PFAS, eaux souterraines, enrichissement g\'eospatial, apprentissage automatique, classification binaire, mode pr\'edictif, transformeur de graphe h\'et\'erog\`ene, valeurs de Shapley, interpr\'etabilit\'e, Californie.

\myCleanStarChapterEnd
